\section{Decorator}
\label{sec:pat-decorator}

\problemstatement{Attach additional responsibilities to an object
dynamically, at run time, without subclassing -- providing a flexible
alternative to inheritance for extending behaviour.}

\begin{patternbox}[When You Reach for It]
The trigger is \emph{``add behaviour at run time in combinations,''} where
a subclass per combination would explode (coffee with milk, with sugar,
with milk-and-sugar-and-cream). When features stack in arbitrary
combinations, wrap rather than subclass.
\end{patternbox}

\subsection*{Understanding the pattern}

A decorator implements the \emph{same interface} as the object it wraps
and holds a pointer to a wrapped component of that interface. It forwards
calls to the inner object and adds its own behaviour before or after. Since
a decorator is itself a component, decorators nest: each layer adds one
feature, and you compose features by stacking wrappers at run time.

\begin{ideabox}[Same Interface, Wrap and Add]
Unlike Adapter (changes the interface) a Decorator keeps the interface
identical and augments the behaviour, so wrapped and unwrapped objects are
interchangeable. Stacking decorators composes features additively without
a single new subclass.
\end{ideabox}

\subsection*{C++ implementation}

\begin{lstlisting}
#include <bits/stdc++.h>
using namespace std;

struct Coffee {                                // component
    virtual double cost() const = 0;
    virtual string desc() const = 0;
    virtual ~Coffee() = default;
};

struct Espresso : Coffee {                     // concrete component
    double cost() const override { return 2.0; }
    string desc() const override { return "Espresso"; }
};

struct CoffeeDecorator : Coffee {              // base decorator
    unique_ptr<Coffee> inner;
    CoffeeDecorator(unique_ptr<Coffee> c) : inner(move(c)) {}
};

struct Milk : CoffeeDecorator {
    Milk(unique_ptr<Coffee> c) : CoffeeDecorator(move(c)) {}
    double cost() const override { return inner->cost() + 0.5; }
    string desc() const override { return inner->desc() + " + Milk"; }
};
struct Sugar : CoffeeDecorator {
    Sugar(unique_ptr<Coffee> c) : CoffeeDecorator(move(c)) {}
    double cost() const override { return inner->cost() + 0.2; }
    string desc() const override { return inner->desc() + " + Sugar"; }
};

int main() {
    unique_ptr<Coffee> order =
        make_unique<Sugar>(make_unique<Milk>(make_unique<Espresso>()));
    cout << order->desc() << " = $" << order->cost() << "\n";  // stacked
    return 0;
}
\end{lstlisting}

\begin{complexitybox}[Trade-offs]
\textbf{Pros.} Add/remove responsibilities at run time; avoids a subclass
per feature-combination; each decorator is a small, single-purpose class.
\textbf{Cons.} Many small objects and layers can be hard to debug; order
of wrapping can matter; a deep stack obscures the concrete object.
\end{complexitybox}

\begin{takeawaybox}
Decorator wraps an object in another of the same interface to add
behaviour, and nests to compose features -- the run-time alternative to a
combinatorial subclass tree. Keyword: ``add responsibilities dynamically.''
It is the pattern behind Java's I/O streams and many middleware chains.
\end{takeawaybox}
